\section{문제 정의와 보장 경계}\label{sec:problem}
입력 $c_t=(G_t,q_t)$에서 $G_t$는 형식이 지정된 게임 스냅샷(장소, 도달 가능 위치, 물체 배치, 소지품, 사실, NPC 지식, 퀘스트 단계)을, $q_t$는 외부에서 저작된 행동 정책·정보 공개·퀘스트 단계 제약을 나타낸다. 검증 술어는 직전 이력 $h_{<t}$를 역산하거나 $q_t$를 변경하지 않으며, 검색 결과나 정서 점수는 쓰기 권한이 없다.

비신뢰 제안 $a_t$를 엄격히 파싱한 뒤, 검증기는 7개의 결정론적 술어를 실행하여 \texttt{ValidationResult}(\allowbreak$\mathrm{valid},E_t,C_{\mathrm{run}}$)를 반환한다:
\begin{equation}
\begin{aligned}
V(c_t,a_t)&=(\mathrm{valid},E_t,C_{\mathrm{run}}),\\
\boldsymbol{v}&=(v_{\mathrm{policy}},v_{\mathrm{pre}},v_{\mathrm{reach}},\\
&\qquad v_{\mathrm{know}},v_{\mathrm{disc}},v_{\mathrm{elig}},v_{\mathrm{stage}})\in\{0,1\}^{7},
\end{aligned}
\end{equation}
여기서 $\mathrm{valid}\equiv\bigwedge_{i=1}^{7}(v_i=1)$이며, $E_t$는 오류 코드·개체·원인·수리 힌트를 기록한다. 정본 상태는 7개 술어를 모두 통과한 경우에만 갱신된다:
\begin{equation}
c_{t+1}=\begin{cases}
T(c_t,a_t), & \bigwedge_{i=1}^{7}(v_i=1),\\
c_t, & \text{otherwise.}
\end{cases}
\label{eq:transition}
\end{equation}

\subsection{스키마 유효성, 권한, 목표 달성의 분리}\label{sec:authorization-sets}
행동 전 상태 $c_t$, 저작 정책, 엄격 스키마를 고정하고 스키마 유효 후보 집합을 $\SchemaSet$라 할 때, 권한과 목표 달성을 다음과 같이 분리한다:
\begin{equation}
\begin{aligned}
\AcceptSet(c_t)&=\{a\in\SchemaSet:V(c_t,a).\mathrm{valid}\},\\
\GoalSet(c_t)&=\{a\in\SchemaSet:\operatorname{Goal}(T(c_t,a))\},\\
\HitSet(c_t)&=\AcceptSet(c_t)\cap\GoalSet(c_t)\subseteq\AcceptSet(c_t)\subseteq\SchemaSet.
\end{aligned}
\end{equation}
금지 사실을 선언한 스키마 유효 후보는 $\AcceptSet$ 밖에 놓이며, 허용된 관찰 행동은 $\AcceptSet$에 속하지만 $\GoalSet$ 밖에 놓일 수 있다. 지식과 공개 권한은 독립 입력이다. 예컨대 NPC는 0단계와 1단계 모두에서 \texttt{keeper\_\allowbreak secret}을 알고 있으나, 공개는 0단계에서 금지되고 1단계에서 허용된다.

\subsection{불변식 I1--I4에 대한 형식 안전성 정리}\label{sec:formal-proofs}
결정론적 전이 규칙(식~\eqref{eq:transition})과 엄격 파서 \texttt{parse\_}\allowbreak\texttt{candidate\_}\allowbreak\texttt{mapping} 아래에서 불변식 \textbf{I1--I4}는 임의의 유한 인코딩 상태 $c_t=(G_t,q_t)$, 수리 예산 $K\in\mathbb{Z}_{\geq 0}$ 및 정상 반환 콜백에 대한 조건부 구조 정리로 성립한다.

\noindent\textbf{정리 1 (I1: 상태 보존 안전성).} \emph{초기 후보가 엄격 파싱을 통과하지 못하거나 $K{+}1$회 모든 시도에서 $V(c_t,a^{(j)}).\mathrm{valid}=\mathrm{false}$이면 $c_{t+1}=c_t$이며 커밋된 효과는 비어 있다.}
\par\noindent\textit{증명.} \texttt{parse\_}\allowbreak\texttt{candidate\_}\allowbreak\texttt{mapping}은 커밋 전에 비정형 제안을 거부하며, 그림~\ref{alg:commit}에서 $\operatorname{Commit}$은 $\mathrm{ok}=\mathrm{true}$일 때만 도달 가능하고 $j=K$에서는 $\operatorname{Fallback}(c_t)$이 식~\eqref{eq:transition}에 의해 $c_t$를 그대로 반환한다. \hfill$\square$

\noindent\textbf{정리 2 (I2: 유한 종료와 상태 비참조 세탁 방지 수리).} \emph{유한 $K{\geq}0$과 정상 반환 콜백 하에서 루프는 $V$를 최대 $K{+}1$회 호출하며, 상태 변경 직전 $V(c_t,a^{(j)})$를 재검증하고, $\rho$는 \texttt{WorldState} 의존 권한 오류 클래스를 변경하지 않는다.}
{\sloppy\par\noindent\textit{증명.} 계수기 $j$는 $0$에서 시작하여 증가 전 $j=K$를 검사하므로 $V$ 호출은 최대 $K{+}1$회로 유계다. 표~\ref{tab:rho-rules}에 따라 $\rho(a,E)$는 \texttt{WorldState}를 읽지 않고 $(a,E)$만 입력받아 \texttt{UNKNOWN\_}\allowbreak\texttt{ACTION\_}\allowbreak\texttt{TYPE}, \texttt{UNREACHABLE\_}\allowbreak\texttt{REQUIRED\_}\allowbreak\texttt{OBJECT}, \texttt{NPC\_}\allowbreak\texttt{KNOWLEDGE\_}\allowbreak\texttt{VIOLATION}, \texttt{NPC\_}\allowbreak\texttt{DIS\-CLOSURE\_}\allowbreak\texttt{KNOW\-LEDGE\_}\allowbreak\texttt{VIO\-LATION}, \texttt{FORBIDDEN\_}\allowbreak\texttt{DISCLOSURE}, \texttt{QUEST\_}\allowbreak\texttt{STAGE\_}\allowbreak\texttt{VIOLATION}, \texttt{QUEST\_}\allowbreak\texttt{STAGE\_}\allowbreak\texttt{REGRESSION}을 항등 사상으로 두고 동일 필드 상충 편집을 제거한다. \hfill$\square$\par}

\noindent\textbf{정리 3 (I3: 정책 유계 상태 변이).} \emph{커밋된 모든 전이는 7개 술어 $v_1\dots v_7$을 만족하며 \texttt{ActionPolicy.}\allowbreak\texttt{allowed\_}\allowbreak\texttt{effects}에 허용된 효과만 반영한다.}
\par\noindent\textit{증명.} 엄격 파서가 12개 \texttt{CANDIDATE\_}\allowbreak\texttt{FIELDS}를 동결하고, $v_{\mathrm{policy}}{=}1$이 허용 집합 포함 여부를 검증하며, $T$는 선언 외 키를 기록하지 않는다. \hfill$\square$

\noindent\textbf{정리 4 (I4: 결정론적 의미 리플레이 및 변조 탐지).} \emph{\texttt{allow\_}\allowbreak\texttt{defaults=}\allowbreak\texttt{false} 파싱과 동결 $(V,T)$ 하에서 정렬된 해시 연쇄 전이 기록을 $c_0$부터 재실행하면 $c_T$가 유일하게 복원되며 단일 단계 변조도 SHA-256 연쇄 불일치로 탐지된다.}
\par\noindent\textit{증명.} $V,T$는 $(c_t,a_t)$의 순수 결정론적 함수이고 각 기록은 $(\mathrm{hash}(c_t),\allowbreak a_t,\allowbreak V(c_t,a_t),\allowbreak\mathrm{hash}(c_{t+1}))$을 해싱하므로 귀납적으로 $\mathrm{hash}(c'_T)=\mathrm{hash}(c_T)$가 일치한다. \hfill$\square$

\section{TRACE-RPG 아키텍처와 알고리즘}\label{sec:arch}
모델과 어댑터가 직접 HTTPS 구조화 디코딩 또는 제한 프롬프트로 JSON 후보 행동을 제안하면, 엄격 파서가 미등록 키와 모호한 타입을 거부하고, 버전 관리 브리지가 관측·제안·검증·수리·대체·커밋을 교환한다(그림~\ref{fig:architecture}, 표~\ref{tab:validators}). 여덟 역할 읽기 전용 온톨로지 뷰(\texttt{World}, \texttt{Actor}, \texttt{Knowledge}, \texttt{Fact}, \texttt{Policy}, \texttt{Quest}, \texttt{Action}, \texttt{Trace})는 쓰기 권한 없이 출처 선언을 실행 검사에 연결한다.

\subsection{다중 홉 Datalog 폐쇄 및 비교 안전 기준선}\label{sec:multihop-baselines}
서사 비밀이 깊이 $d$까지의 다단계 사실 함의 또는 역할 상속 연쇄를 거쳐 파생될 때, 평면 1홉 검사만으로는 간접 누출을 막지 못한다. 핵심 서브패키지 \texttt{nesy\_\allowbreak game.\allowbreak multihop}은 불변 사양 \texttt{Multihop\allowbreak Ontology\allowbreak Spec}(\texttt{role\_\allowbreak parents}, \texttt{fact\_\allowbreak implications}, \texttt{forbidden\_\allowbreak by\_\allowbreak role}, \texttt{min\_\allowbreak stage\_\allowbreak by\_\allowbreak fact})에 대해 \texttt{compute\_\allowbreak datalog\_\allowbreak closure}를 호출하여 상향 준단순(semi-naive) 고정점 폐쇄를 수행한다:
\begin{equation}
\mathcal{F}^*(S)=\operatorname{lfp}\Big(X\mapsto S\cup\bigcup_{f\in X}\texttt{fact\_implications}(f)\Big),
\end{equation}
그리고 역할 상속 폐쇄 $\mathcal{R}^*(u)$의 금지 규칙과 상태 금지 집합을 합산한 뒤 \texttt{validate\_\allowbreak candidate\_\allowbreak multihop}에서 $v_{\mathrm{know}}, v_{\mathrm{disc}}, v_{\mathrm{stage}}$를 판정한다. 외부 신경망 강화학습 훈련이나 별도 서버 런타임 없이 동일 \texttt{WorldState} 계약 위에서 안전 패러다임을 공정 대비하기 위해 패키지 내 두 통제 대조 프록시 구현을 탑재하였다: (i)~\textbf{Safe-RL Action Masking}(\texttt{evaluate\_\allowbreak action\_\allowbreak mask\_\allowbreak shield}~\cite{alshiekh2018shield})은 행동 유형·도달 가능성·1홉 선행조건 마스크를 강제하나 이행적 역할·사실 폐쇄를 계산하지 않는다; (ii)~\textbf{NeMo 스타일 Colang 출력 레일}(\texttt{evaluate\_\allowbreak dialog\_\allowbreak rail\_\allowbreak filter}~\cite{rebedea2023nemo})은 발화 표면 문자열과 직접 1홉 대화 정책만 검사한다.

\begin{table}[!tbp]
\caption{인코딩된 7술어 검증기 경계}
\label{tab:validators}
\centering
\footnotesize
\begin{tabularx}{\columnwidth}{@{}>{\raggedright\arraybackslash}p{0.24\columnwidth}>{\raggedright\arraybackslash}X>{\raggedright\arraybackslash}p{0.16\columnwidth}@{}}
\toprule
\textbf{술어} & \textbf{거부 조건} & \textbf{실패 시} \\
\midrule
행동 정책 & 미등록 행동/유형, 필수 조건·효과 누락, 비허용 효과 & 상태 불변 \\
선행조건 & 부재 또는 거짓인 필수 조건 & 상태 불변 \\
도달 가능성 & 접근 불가능한 대상 물체 & 상태 불변 \\
NPC 지식 & 행위자 지식에 없는 사실 사용·공개 & 상태 불변 \\
정보 공개 & 선언된 금지 사실 및 Datalog 파생 금지 사실 & 상태 불변 \\
퀘스트 단계 & 현재 단계보다 높은 요구/파생 최소 단계 & 상태 불변 \\
단계 단조성 & 역행하는 퀘스트 단계 전이 & 상태 불변 \\
\bottomrule
\end{tabularx}
\end{table}

\begin{figure}[!tbp]
\centering
\footnotesize
\begin{tabular}{@{}rl@{}}
1: & $a\leftarrow\operatorname{Propose}(c)$; $j\leftarrow0$ \\
2: & \textbf{repeat} \\
3: & \quad $(\mathrm{ok},E)\leftarrow V(c,a)$; $\operatorname{Record}(a,E)$ \\
4: & \quad \textbf{if} $\mathrm{ok}$ \textbf{then assert} $V(c,a).\mathrm{ok}$; \textbf{return} $\operatorname{Commit}(c,a)$ \\
5: & \quad \textbf{if} $j=K$ \textbf{then return} $\operatorname{Fallback}(c)$ \\
6: & \quad $a\leftarrow\operatorname{Repair}(c,a,E)$; $j\leftarrow j+1$ \\
7: & \textbf{until terminal}
\end{tabular}
\caption{검증--수리--커밋 루프는 최대 $K+1$회 시도를 기록하고 커밋 직전에 재검증한다.}
\Description{최대 K+1회 시도를 기록하고 커밋 직전 재검증하는 검증-수리-커밋 루프 의사코드.}
\label{alg:commit}
\end{figure}

유도 수리 연산자 $\rho(a,E)=a\oplus\bigcup_{e\in E}\delta(e)$(그림~\ref{alg:commit}, 표~\ref{tab:rho-rules})와 \texttt{execute\_\allowbreak with\_\allowbreak multihop\_\allowbreak repair}는 \texttt{WorldState}를 전혀 읽지 않고 후보 $a$와 구조 오류 $E$만 입력받아 우회 세탁(state laundering)을 차단한다. 전이 기록은 이전 상태, 후보, 검증 결과, 재계산된 최종 상태에 대해 비키 SHA-256 해시를 연쇄 기록해 결정론적 리플레이를 지원한다.

\begin{table}[!tbp]
\caption{유도 수리 연산자 $\rho(a,E)$의 오류 유형별 편집 규칙}
\label{tab:rho-rules}
\centering
\footnotesize
\begin{tabularx}{\columnwidth}{@{}>{\raggedright\arraybackslash}X>{\raggedright\arraybackslash}X@{}}
\toprule
\textbf{검증기 오류 유형} & \textbf{후보 $a$에 대한 편집 $\delta(e)$} \\
\midrule
정책 선행조건 / 효과 누락 & 명시된 술어 또는 효과 삽입 \\
정책 효과 위반 & 비허용 효과 제거 \\
불만족 선행조건 & 명시된 술어 제거 \\
비인가 단계 변경 & \texttt{quest\_stage\_effect} 초기화 \\
미등록 행동; 도달·지식·공개·단계 오류 & 무연산(no-op, 상태 세탁 방지) \\
\bottomrule
\end{tabularx}
\end{table}
